% Partie : set-theory
% Chapitre : ordinals
% Section : introduction

\documentclass[../../../include/open-logic-section]{subfiles}

\begin{document}

\olfileid[fr]{sth}{ordinals}{intro}

\olsection{Introduction}

Dans le \olref[z][]{chap}, nous avons postulé qu'il existe une étape
infinie de la hiérarchie, sous la forme de \stagesinf{} (voir aussi
notre axiome de l'infini). Mais, étant donné \stagessucc{}, nous ne
pouvons pas nous arrêter à cette étape infinie : il faut poursuivre.
À l'étape qui suit la première étape infinie, nous formons donc toutes
les collections possibles d'ensembles qui étaient disponibles à la
première étape infinie ; puis nous recommençons, encore et encore,
\dots

Ce faisant, nous avons implicitement commencé à invoquer une notion
«~intuitive~» de nombre selon laquelle il peut exister des nombres
\emph{après} tous les nombres naturels. Il s'agit, plus précisément,
de la notion d'\emph{ordinal transfini}. Ce chapitre vise à rendre
cette idée plus rigoureuse. Nous étudierons la notion générale
d'ordinal, puis nous définirons explicitement des ensembles pour
représenter les ordinaux.

\end{document}
